% dynamics/Vienna_reinterpretation.tex

\chapter{Vienna Reinterpretation and Dynamic State Shaping}

\label{chap:vienna-reinterpretation}

% ============================================================
\section*{Motivation}
% ============================================================

Classical transition theory usually interprets transition curves as
special geometric constructions.

Within AIM, however, transition structures are interpreted primarily
through their generated runtime behaviour.

This changes the interpretation of Vienna-style transition concepts
fundamentally.

Within this framework, Vienna is not interpreted as:
\begin{itemize}
\item a specific curve formula,
\item a historical construction rule,
\item or a particular interpolation method,
\end{itemize}
but rather as:
\begin{quote}
a strategy for shaping operational runtime states.
\end{quote}

\begin{principle}[Vienna reinterpretation principle]
\label{princ:vienna-reinterpretation}
Within AIM, Vienna-type transitions are interpreted as dynamic
state-shaping strategies operating in curvature-runtime space.
\end{principle}

\begin{principle}[No special-curve principle]
\label{princ:vienna-no-special-curve}
Vienna reinterpretation does not elevate one special curve formula.
It interprets transition structure by its generated operational runtime
behaviour.
\end{principle}

% ============================================================
\section{From Curves to Runtime States}
% ============================================================

Classical railway alignment theory focuses primarily on:
\[
\gamma(s).
\]

Within AIM, however, the operationally relevant quantity is:
\[
\mathrm{pose3}_{\mathrm{real}}(s),
\]
together with the generated vehicle-space trajectory.

Thus, transition interpretation shifts:
\begin{itemize}
\item from coordinate geometry,
\item toward runtime operational state evolution.
\end{itemize}

The central engineering question therefore becomes:
\begin{quote}
Which operational runtime trajectory is generated?
\end{quote}
rather than:
\begin{quote}
Which geometric curve formula is used?
\end{quote}

% ============================================================
\section{Vienna as Runtime-State Strategy}
% ============================================================

\begin{definition}[Vienna-type state shaping]
\label{def:vienna-state-shaping}
A Vienna-type transition strategy is a curvature-evolution strategy
designed to shape operational runtime-state trajectories.
\end{definition}

Relevant runtime quantities include:
\begin{itemize}
\item effective acceleration,
\item jerk,
\item roll evolution,
\item operational balance,
\item and center-of-gravity motion.
\end{itemize}

Thus, Vienna interpretation operates fundamentally in:
\begin{itemize}
\item curvature space,
\item runtime space,
\item and operational state space.
\end{itemize}

The term \emph{Vienna} is therefore used here as an interpretive
strategy, not as a commitment to one fixed analytical transition law.

% ============================================================
\section{Half-Wave Interpretation}
% ============================================================

\begin{definition}[Half-wave interpretation]
\label{def:half-wave-interpretation}
Half-wave interpretation denotes a modular description of transition
evolution through normalized curvature segments.
\end{definition}

Within AIM, many transition structures may be interpreted
compositionally as:
\[
\text{halfWave}_1
\rightarrow
\text{core}
\rightarrow
\text{halfWave}_2.
\]

This decomposition is a modelling aid.
It should not be interpreted as a new special-curve dogma.

The half-wave interpretation allows:
\begin{itemize}
\item modular runtime shaping,
\item asymmetric runtime behaviour,
\item spectral interpretation,
\item and realization-aware transition design.
\end{itemize}

The half-wave structure is relevant only insofar as it affects generated
runtime behaviour.

% ============================================================
\section{Dynamic State Shaping}
% ============================================================

\begin{definition}[Dynamic state shaping]
\label{def:dynamic-state-shaping}
Dynamic state shaping denotes controlled generation of operational
runtime trajectories through curvature evolution.
\end{definition}

Within AIM, transition geometry is interpreted fundamentally as:
\begin{quote}
runtime-state shaping through curvature control.
\end{quote}

Relevant shaped quantities include:
\begin{itemize}
\item effective acceleration,
\item jerk,
\item cant deficiency,
\item roll evolution,
\item and center-of-gravity trajectories.
\end{itemize}

% ============================================================
\section{Center-of-Gravity Interpretation}
% ============================================================

Within Vienna reinterpretation, the operationally relevant object is
not merely the rail geometry itself.

Instead, emphasis shifts toward:
\[
x_{\mathrm{CG}}(s),
\]
the runtime trajectory of the vehicle center of gravity.

Different transition families may produce:
\begin{itemize}
\item similar geometric centerlines,
\item but fundamentally different center-of-gravity trajectories.
\end{itemize}

Thus, Vienna reinterpretation naturally connects transition theory with:
\begin{itemize}
\item Schwerpunkt-Trassierung,
\item operational balance,
\item and runtime comfort interpretation.
\end{itemize}

% ============================================================
\section{Spectral Interpretation}
% ============================================================

Transition laws may be interpreted spectrally through their curvature
content.

Low-frequency curvature behaviour primarily influences:
\begin{itemize}
\item global operational balance,
\item smooth acceleration evolution,
\item and large-scale runtime behaviour.
\end{itemize}

High-frequency curvature behaviour tends to influence:
\begin{itemize}
\item realization sensitivity,
\item machine smoothing,
\item dynamic roughness,
\item measurement sensitivity,
\item and operational excitation.
\end{itemize}

\begin{proposition}[Spectral smoothing principle]
\label{prop:spectral-smoothing}
Operationally smooth runtime trajectories correspond to spectrally
controlled curvature evolution.
\end{proposition}

This interpretation shifts transition evaluation away from geometric
appearance and toward spectral runtime behaviour.

% ============================================================
\section{Realization Filtering}
% ============================================================

Realization acts approximately as a low-pass filter on curvature space.

Thus, highly oscillatory transition structures may collapse into similar
realized operational states after construction and maintenance.

Within AIM, Vienna reinterpretation therefore includes:
\begin{itemize}
\item realization filtering,
\item machine smoothing,
\item measurement reconstruction,
\item and runtime operational robustness.
\end{itemize}

A transition law is operationally meaningful only if its intended
runtime behaviour survives realization.

\begin{principle}[Realization-survival principle]
\label{princ:realization-survival}
A transition strategy is operationally relevant only to the extent that
its intended runtime behaviour survives realization filtering.
\end{principle}

A possible Vienna-type hypothesis is therefore that certain transition
strategies may preserve their intended operational behaviour more robustly
under realization filtering than others.

% ============================================================
\section{Runtime Smoothness}
% ============================================================

\begin{definition}[Runtime smoothness]
\label{def:runtime-smoothness}
Runtime smoothness denotes smooth operational evolution of runtime
states generated through curvature and cant evolution.
\end{definition}

Runtime smoothness includes:
\begin{itemize}
\item smooth acceleration evolution,
\item smooth roll evolution,
\item smooth operational balance,
\item and smooth center-of-gravity trajectories.
\end{itemize}

Within AIM, runtime smoothness is operationally more relevant than
purely geometric smoothness.

Spectral smoothing provides one possible interpretation of runtime
smoothness in curvature space.

% ============================================================
\section{Vienna and Curvature Space}
% ============================================================

Vienna reinterpretation operates fundamentally within curvature space.

Different curvature structures generate:
\begin{itemize}
\item different operational balance trajectories,
\item different roll behaviour,
\item different jerk evolution,
\item different spectral content,
\item and different realization behaviour.
\end{itemize}

Thus, transition quality cannot be characterized sufficiently through
position geometry alone.

% ============================================================
\section{Vienna and Runtime Dynamics}
% ============================================================

Vienna reinterpretation is fundamentally a runtime-dynamics concept.

It couples:
\begin{itemize}
\item curvature evolution,
\item runtime-state evolution,
\item operational balance,
\item spectral smoothing,
\item realization filtering,
\item and vehicle-space interpretation.
\end{itemize}

The transition curve therefore acts as:
\begin{quote}
a runtime-state generator.
\end{quote}

% ============================================================
\section{Operational Interpretation}
% ============================================================

Within AIM, the operational railway problem is interpreted as:
\begin{quote}
generation of admissible runtime operational trajectories.
\end{quote}

This transforms railway alignment theory from:
\begin{itemize}
\item geometric interpolation,
\item toward operational runtime-state design.
\end{itemize}

Vienna reinterpretation therefore acts as a bridge between:
\begin{itemize}
\item transition theory,
\item runtime dynamics,
\item realization,
\item and operational optimization.
\end{itemize}

% ============================================================
\section{Towards Runtime-State Optimization}
% ============================================================

The reinterpretation developed here prepares the transition toward:
\begin{itemize}
\item realization-aware optimization,
\item operational admissibility optimization,
\item runtime-state shaping,
\item and Schwerpunkt-based alignment design.
\end{itemize}

Optimization therefore acts not merely on:
\[
\gamma(s),
\]
but on:
\[
\mathrm{pose3}_{\mathrm{real}}(s),
\]
and the generated operational vehicle-state trajectories.

% ============================================================
\section{Canonical Interpretation}
% ============================================================

\begin{proposition}
Within AIM, Vienna-type transitions are interpreted as runtime-state
shaping strategies rather than as isolated geometric curve formulas.
\end{proposition}

\begin{proposition}
Operational transition quality depends fundamentally on generated
runtime-state trajectories.
\end{proposition}

\begin{proposition}
Realization filtering and spectral behaviour are central components of
transition interpretation.
\end{proposition}

\begin{proposition}
The operationally relevant object is the realized runtime trajectory of
vehicle states rather than the abstract geometric curve alone.
\end{proposition}

\begin{proposition}
The half-wave interpretation is a modular modelling aid, not a
special-curve doctrine.
\end{proposition}

% ============================================================
\section{Connection to AIM}
% ============================================================

\begin{summary}
Vienna reinterpretation transforms transition theory into a theory of
runtime-state shaping.

Within AIM:
\begin{itemize}
\item transitions generate operational runtime trajectories,
\item half-wave structures provide modular interpretation rather than
special-curve dogma,
\item realization filtering determines operational robustness,
\item spectral smoothing influences comfort and admissibility,
\item and center-of-gravity trajectories become primary operational
objects.
\end{itemize}

This establishes the conceptual bridge between:
\begin{itemize}
\item curvature space,
\item runtime dynamics,
\item realization-aware interpretation,
\item Schwerpunkt-Trassierung,
\item and operational state optimization.
\end{itemize}
\end{summary}
